\subsection{Coefficient conditioning and identifiability conditional on the selected support}
\label{subsec:d3d-coefficient-conditioning}

Conditional on the frozen seven-coordinate support of the canonical DIII-D
reconstruction package (\texttt{D3D-SIR-62-ALIGNED-V1}), we assessed whether
discharge-specific coefficients are numerically identifiable and whether
cross-discharge coefficient variation exceeds within-discharge estimation
uncertainty. Feature ordering is tracked separately as source-matrix order and
manuscript display order.

Least-squares reproduction of the canonical discharge-specific fits recovers
pooled RMSE $0.05758467$. The cohort-mean coefficient
vector recovers pooled RMSE $0.41252383$.
All 62 design matrices have numerical rank~7 at relative singular-value
tolerances from $10^{-12}$ through $10^{-6}$, with median spectral
condition number $41.32$.

Original relative truncation thresholds $\tau\le 10^{-6}$ removed no singular
directions: the smallest normalized singular values satisfy
$\sigma_7/\sigma_1\gtrsim 0.0140$.
Meaningful fixed-rank truncation shows that removing the weakest direction
produces a median relative coefficient-vector change of $0.799$ and a
median absolute RMSE change of $0.031$. Thus reconstruction can remain
comparatively stable while coefficients move substantially along weakly resolved
directions. A ridge path on standardized features yields the same qualitative
conclusion under continuous shrinkage.

The original random-effects column labeled REML implemented profile maximum
likelihood. Corrected analysis uses heteroscedastic restricted maximum
likelihood with known discharge-specific block-bootstrap variances, profile
intervals, short/long block sensitivity, and leave-one-discharge-out checks.
Under these criteria, 5 coefficients are robustly resolved
(dot kappa, D_betaN W_dia, D_betaN kappa, dot beta_N, D_betaN l_i), while 2 remain dominated by within-discharge
uncertainty (D_kappa W_dia, q95 / kappa).

A multivariate analysis of $\Sigma_B-\overline{\Sigma}_W$ originally counted
5 positive eigenvalues.
After a 5000-replicate discharge bootstrap and a
matched no-between-discharge-variation null
(5000 replicates), 2 directions remain robustly
resolved. Positive point estimates alone are not treated as identified
directions.

Overall, discharge-specific reconstruction is stable, but coefficient
identifiability is mixed: several coefficients and a lower-dimensional subset of
coefficient-space directions are resolved after uncertainty correction, while
others are not. This analysis does not test uniqueness of the selected support
and does not establish mechanistic or causal interpretation.
